% Part: first-order-logic
% Chapter: tableaux
% Section: identity

\documentclass[../../../include/open-logic-section]{subfiles}

\begin{document}

\olfileid{fol}{tab}{ide}

\olsection{\usetoken{S}{identity} સાથેના \usetoken{P}{tableau}}

!!{identity} ધરાવતા !!^{tableau}s માટે વધારાના અનુમાન-નિયમો જરૂરી છે.
$\eq$ માટેના નિયમો આ છે ($t$, $t_1$ અને $t_2$ બંધ પદો છે):

\begin{defish}
\AxiomC{}
\RightLabel{$\eq$}
\UnaryInfC{\sFmla{\True}{\eq[t][t]}}
\DisplayProof
\hfill
\AxiomC{\sFmla{\True}{\eq[t_1][t_2]}}
\noLine
\UnaryInfC{\sFmla{\True}{!A(t_1)}}
\RightLabel{$\TRule{\True}{\eq}$}
\UnaryInfC{\sFmla{\True}{!A(t_2)}}
\DisplayProof
\hfill
\AxiomC{\sFmla{\True}{\eq[t_1][t_2]}}
\noLine
\UnaryInfC{\sFmla{\False}{!A(t_1)}}
\RightLabel{$\TRule{\False}{\eq}$}
\UnaryInfC{\sFmla{\False}{!A(t_2)}}
\DisplayProof
\end{defish}
નોંધો કે બીજા બધા નિયમોથી વિપરીત, $\TRule{\True}{\eq}$ અને
$\TRule{\False}{\eq}$ માટે શાખા પર \emph{બે} ચિહ્નિત !!{formula}s
પહેલેથી આવવાં જરૂરી છે: $\sFmla{\True}{\eq[t_1][t_2]}$ અને
$\sFmla{S}{!A(t_1)}$ બંને.

\begin{ex}
જો $s$ અને $t$ બંધ પદો હોય, તો $\eq[s][t], !A(s) \Proves !A(t)$:
\begin{oltableau}
  [\sFmla{\False}{\formula{A}(t)}, just = \TAss
    [\sFmla{\True}{\eq[s][t]}, just = \TAss
      [\sFmla{\True}{\formula{A}(s)}, just = \TAss
        [\sFmla{\True}{\formula{A}(t)}, just={\TRule{\True}{\eq}[2, 3]}, close]
      ]
    ]
  ]
\end{oltableau}
આ અભિન્ન વસ્તુઓની પ્રતિસ્થાપનીયતાના સિદ્ધાંત તરીકે, અથવા લાઇબ્નિત્સના
નિયમ તરીકે, પરિચિત હોઈ શકે છે.

!!^{tableau}s સાબિત કરે છે કે $\eq$ સંમિત છે, એટલે કે
$\eq[s_1][s_2] \Proves \eq[s_2][s_1]$:
\begin{oltableau}
  [\sFmla{\False}{\eq[s_2][s_1]}, just = \TAss
    [\sFmla{\True}{\eq[s_1][s_2]}, just = \TAss
      [\sFmla{\True}{\eq[s_1][s_1]}, just = {$\eq$}
        [\sFmla{\True}{\eq[s_2][s_1]}, just = {\TRule{\True}{\eq}[2, 3]}, close]
      ]
    ]
  ]
\end{oltableau}
અહીં પંક્તિ $2$ એ $\TRule{\True}{\eq}$ માટે જરૂરી પહેલું
!!{formula}~$\sFmla{\True}{\eq[s_1][s_2]}$ છે. પંક્તિ~$3$ એ
$\sFmla{\True}{!A(s_1)}$ સ્વરૂપનું બીજું આવશ્યક સૂત્ર છે---$!A(x)$ને
$\eq[x][s_1]$ તરીકે વિચારો; ત્યારે $!A(s_1)$ એ $\eq[s_1][s_1]$ અને
$!A(s_2)$ એ $\eq[s_2][s_1]$ છે.
\sourcecorrection{OLTAB-011}{મૂળની \texttt{identity.tex}, પંક્તિ~69માં
નિયમ માટેનું બીજું આવશ્યક સૂત્ર $\sFmla{\True}{!A(s_2)}$ લખાયું છે.
અહીં $t_1=s_1$, $t_2=s_2$ અને તરત પછીનું
$!A(x)=\eq[x][s_1]$ હોવાથી આધારવાક્ય $!A(s_1)$ અને ફલિત $!A(s_2)$ છે;
તેથી આધારવાક્યનો સૂચકાંક પુનઃસ્થાપિત કર્યો છે.}

તેઓ એ પણ સાબિત કરે છે કે $\eq$ પરંપરિત છે, એટલે કે
$\eq[s_1][s_2], \eq[s_2][s_3] \Proves \eq[s_1][s_3]$:
\begin{oltableau}
  [\sFmla{\False}{\eq[s_1][s_3]}, just = \TAss
    [\sFmla{\True}{\eq[s_1][s_2]}, just = \TAss
      [\sFmla{\True}{\eq[s_2][s_3]}, just = \TAss
        [\sFmla{\True}{\eq[s_1][s_3]}, just = {\TRule{\True}{\eq}[3, 2]}, close]
      ]
    ]
  ]
\end{oltableau}
આ !!{tableau}માં $\TRule{\True}{\eq}$ માટે જરૂરી પહેલું !!{formula}
પંક્તિ~$3$ પરનું $\sFmla{\True}{\eq[s_2][s_3]}$ છે ($s_2$ એ~$t_1$ની
અને $s_3$ એ~$t_2$ની ભૂમિકા ભજવે છે). $\sFmla{\True}{!A(s_2)}$
સ્વરૂપનું બીજું આવશ્યક સૂત્ર પંક્તિ~$2$ છે. અહીં $!A(x)$ને
$\eq[s_1][x]$ તરીકે વિચારો; તેથી $!A(s_2)$ એ $\eq[s_1][s_2]$ (એટલે
કે પંક્તિ~$2$) અને $!A(s_3)$ એ ફલિતમાંનું !!{formula}
~$\eq[s_1][s_3]$ બને છે.
\sourcecorrection{OLTAB-010}{મૂળની \texttt{identity.tex}, પંક્તિ~89માં
$!A(s_2)$ને $\eq[t_1][t_2]$ કહ્યું છે. અહીં નિર્ધારિત
$!A(x)=\eq[s_1][x]$ અને પંક્તિ~2 મુજબ તે $\eq[s_1][s_2]$ છે; એ જ
પુનઃસ્થાપિત કર્યું છે.}
\end{ex}

\begin{prob}
નીચેના માટે બંધ !!{tableau}s આપો:
\begin{enumerate}
\item $\sFmla{\False}{\lforall[x][\lforall[y][((x = y \land !A(x))
      \lif !A(y))]]}$
\item $\sFmla{\False}{\lexists[x][(!A(x) \land
    \lforall[y][(!A(y) \lif y = x)])]}$,\\
  $\sFmla{\True}{\lexists[x][!A(x)] \land
    \lforall[y][\lforall[z][((!A(y) \land !A(z)) \lif y = z)]]}$
\end{enumerate}
\end{prob}

\end{document}
